\documentclass[11pt]{article}
\usepackage[margin=1in]{geometry}
\usepackage{amsmath,amssymb,amsthm}
\usepackage{booktabs}
\usepackage{hyperref}
\hypersetup{hidelinks}

\newtheorem{theorem}{Theorem}
\newtheorem{lemma}[theorem]{Lemma}
\newtheorem{corollary}[theorem]{Corollary}
\newtheorem{proposition}[theorem]{Proposition}
\theoremstyle{definition}
\newtheorem{definition}{Definition}
\theoremstyle{remark}
\newtheorem{remark}{Remark}

\setlength{\parskip}{4pt plus 1pt minus 1pt}
\setlength{\parindent}{0pt}

\title{Cancer as Escalation Chain Severance\\Under Bounded Context}
\author{Patrick D.\ McCarthy}
\date{}

\begin{document}
\maketitle

\begin{abstract}
We apply the bounded-context framework of Papers~1--4 to
multicellular organization, proposing that cancer is a predictable
failure mode of any multi-level escalation hierarchy. In a
multicellular organism, individual cells participate in an escalation
chain -- cell $\to$ tissue $\to$ organ $\to$ organism -- where
each level extends the effective bounded context
$C_n^{\mathrm{eff}}$ of the level below via intercellular signaling.
We argue that severing the chain at level $k$ forces the cell back
to $C_n^{(k)}$-level optimization, where the action space available under
cell-level bounded context converges on the recognized hallmarks of
cancer: aerobic glycolysis, unconstrained proliferation, immune
evasion, and motility. This convergence is not atavistic reversion
to ancient programs but \emph{structurally predicted}: the same
constraint ceiling that governed unicellular life now governs the
severed cell. The framework offers structural explanations for three
phenomena that existing theories address incompletely: (1)~the
velocity of malignant transformation (structural failure, not
stochastic accumulation), (2)~the convergent phenotype across tissue
types (same constraint, same necessary structure), and (3)~the
therapeutic resistance of established tumors (locally rational
optimization under cell-level $C_n$). Clinical oncology is already
trending toward the therapeutic strategies this framework predicts
-- differentiation therapy, immunotherapy, and microenvironment
normalization -- but largely without a unified theoretical
justification. We propose that the bounded-context framework
provides one, and offer testable predictions that distinguish this
account from the somatic mutation theory, the atavistic theory, and
the tissue organization field theory.
\end{abstract}

% ------------------------------------------------------------------
\section{Introduction}
\label{sec:intro}
% ------------------------------------------------------------------

Cancer remains the central unsolved problem of multicellular biology.
Three major theoretical frameworks compete to explain it, each
capturing part of the phenomenon while missing the rest.

The \emph{somatic mutation theory} (SMT) holds that cancer arises from
accumulated genetic mutations that progressively disable growth
controls \cite{Vogelstein2013, Nowell1976}. SMT explains the
stochastic component of cancer incidence and the role of specific
oncogenes and tumor suppressors. It does not explain the \emph{velocity}
of malignant transformation: why cancers from different tissues, with
different mutational profiles, converge on the same phenotypic
hallmarks \cite{Hanahan2011} far faster than random mutation would
predict.

The \emph{atavistic theory} holds that cancer cells reactivate ancient
genetic programs from the unicellular era, effectively reverting to a
pre-multicellular phenotype \cite{Davies2011, Trigos2017}. The
phylostratigraphic evidence is striking: genes upregulated in cancer
are disproportionately ancient, dating to unicellular ancestors
\cite{Trigos2017}. But the theory lacks a mechanism: \emph{why} would
a modern cell access billion-year-old programs, and \emph{how} are
those programs still coherent after extensive evolutionary modification?

The \emph{tissue organization field theory} (TOFT) holds that cancer
is not a disease of individual cells but of tissue organization:
the breakdown of intercellular signaling and tissue architecture
enables default proliferative behavior \cite{Sonnenschein2000}. TOFT
correctly identifies the organizational level at which cancer operates
but lacks a formal account of \emph{why} organizational breakdown
produces the specific phenotypes it does.

This paper proposes that all three theories describe aspects of a
single structural phenomenon: \textbf{escalation chain severance}.
We model multicellular organization as a multi-level escalation
hierarchy where each level extends the effective \emph{bounded
context} of the level below. Severing the chain is a structural
failure that forces cell-level optimization. The convergent
phenotype, on this account, is not inherited from ancestors but
structurally predicted under cell-level bounded context. The same
constraint that shaped unicellular life now shapes the severed cell.

The theoretical foundation draws on the bounded-context framework
\cite{McCarthy2026a, McCarthy2026b, McCarthy2026c, McCarthy2026d},
which derives from a single constraint---bounded active context---the
necessity of graph-structured knowledge, escalation hierarchies,
knowledge-action coupling, and adaptive descent. This paper applies
those results to cancer biology. However, the paper is designed to be
self-contained: Section~\ref{sec:framework} develops the framework
from first principles in biological terms, and the central
predictions (Section~\ref{sec:predictions}) are testable
independently of whether one accepts the formal proofs in
\cite{McCarthy2026a}--\cite{McCarthy2026d}.

\textbf{Scope and epistemology.} This paper is a theoretical
application to cancer biology. We do not claim to prove anything
about cancer. We claim that the framework generates a coherent
structural account that (a)~unifies three competing theories,
(b)~aligns with the direction clinical oncology is already moving,
(c)~produces testable predictions that distinguish it from existing
accounts, and (d)~makes a novel central prediction---coupling-channel
specificity (Section~\ref{sec:predictions})---that no existing cancer
theory generates. The value of the paper is the structural reasoning
and the predictions, not empirical validation, which is the province
of cancer biologists and clinicians.

% ------------------------------------------------------------------
\section{Formal Framework}
\label{sec:framework}
% ------------------------------------------------------------------

We develop the framework from three premises that are independently
motivated by biology, then apply them to multi-level biological
organization. The formal derivations of these premises from
information-theoretic first principles appear in
\cite{McCarthy2026a, McCarthy2026b, McCarthy2026c, McCarthy2026d};
here we state them in biological terms.

\textbf{Premise 1: Bounded active context ($C_n$).}
Any information-processing entity---cell, tissue, organism---has a
hard limit on the number of regulatory states it can simultaneously
maintain for active inference. A cell can monitor a finite number of
signaling pathways concurrently. An organism's working memory holds
$4 \pm 1$ items \cite{Cowan2001}. We denote this limit $C_n$.

\textbf{Premise 2: Retention under displacement.}
When $C_n$ is full and new information arrives, something must be
displaced. Under survival pressure, the retention criterion is
asymmetric \cite{McCarthy2026c}: retain a regulatory state when its
survival benefit is uncertain; displace it only when a new state
offers strictly higher marginal survival value. The cost of losing a
survival-relevant state (death) vastly exceeds the cost of retaining
a useless one (a wasted context slot). This drives displacement
toward the highest-value regulatory states available to the entity.

\textbf{Premise 3: Adaptive descent.}
An entity with $\gamma > 0$ (nonzero probability of engaging its
adaptive mechanism in response to signals) and $\eta_M > 0$
(positive learning efficiency) reduces its uncertainty about
survival-relevant states over time \cite{McCarthy2026d}. The entity
performs local descent on uncertainty: each adaptive response, on
average, improves its model of the world. This is not an assumption
about intelligence; it is a minimal condition for persistence. Entities
with $\gamma = 0$ (no adaptive response) or $\eta_M = 0$ (no
learning) are eliminated under positive novelty rate
\cite{McCarthy2026d}.

From these three premises, we model multi-level biological
organization.

\begin{definition}[Biological Entity at Level $k$]
\label{def:bio-entity}
A biological entity at organizational level $k$ is
$E_k = (K_k, A_k, \sigma_k, \gamma_k)$ with bounded active context
$C_n^{(k)}$, where:
\begin{itemize}
\item $K_k$ is the knowledge base at level $k$ (the regulatory state
   accessible for inference: gene expression profiles, signaling
   molecule concentrations, tissue-level morphogen gradients).
\item $A_k$ is the action space at level $k$ (the executable responses
   --- protein synthesis, cell division, apoptosis, migration,
   differentiation signals).
\item $\sigma_k$ is the sensing function at level $k$ (receptors,
   gap junctions, paracrine signaling).
\item $\gamma_k$ is the agency at level $k$ (the probability that
   $E_k$ engages its adaptive mechanism in response to signals).
\end{itemize}
The organizational levels in a metazoan are:
\begin{align*}
k = 0 &: \text{cell} \\
k = 1 &: \text{tissue} \\
k = 2 &: \text{organ} \\
k = 3 &: \text{organism}
\end{align*}
\end{definition}

\begin{definition}[Escalation Chain]
\label{def:chain}
An \emph{escalation chain} is a sequence of coupled entities
$E_0, E_1, \ldots, E_m$ where level $k+1$ receives escalation
signals from level $k$ when $E_k$ encounters states exceeding its
local capacity:
\[
U_k(w, K_k) > U_{\mathrm{threshold}}^{(k)}
\;\Rightarrow\;
\text{escalate to } E_{k+1}
\]
This is the biological instantiation of escalation dominance
\cite{McCarthy2026b}: higher levels engage only when lower levels
fail, achieving $O(1)$ steady-state context cost.
\end{definition}

\begin{definition}[Effective Context Extension]
\label{def:extension}
When $E_k$ is coupled to $E_{k+1}$ via an intact escalation chain,
the \emph{effective bounded context} of $E_k$ is:
\[
C_n^{\mathrm{eff}}(E_k) = C_n^{(k)} + \sum_{j=k+1}^{m}
\lambda_j \cdot C_n^{(j)}
\]
where $\lambda_j \in [0, 1]$ is the \emph{aggregate} coupling strength
between levels, summarizing the fidelity and bandwidth of intercellular
signaling across all channels at level $j$
(Section~\ref{sec:coupling-biology} decomposes $\lambda_j$ into
per-channel components). $\lambda_j = 1$ represents perfect coupling
(all of $E_j$'s context is available to $E_k$ via signaling);
$\lambda_j = 0$ represents complete decoupling.

The key property: $C_n^{\mathrm{eff}} > C_n^{(k)}$. The cell
``borrows'' context from higher organizational levels. A liver cell
does not need to independently represent the organism's metabolic
state; it receives that information via hormonal and neural
signaling from the organ and organism levels.
\end{definition}

\begin{definition}[Chain Severance]
\label{def:severance}
\emph{Chain severance at level $k$} occurs when the coupling
$\lambda_{k+1} \to 0$: the entity $E_k$ loses its connection to all
higher organizational levels. Post-severance:
\[
C_n^{\mathrm{eff}}(E_k) = C_n^{(k)}
\]
The cell's effective context collapses to its intrinsic cell-level
bound.
\end{definition}

\subsection{Biological Instantiation of the Coupling Coefficients}
\label{sec:coupling-biology}

The abstract coupling coefficients $\lambda_j$ correspond to
specific, measurable intercellular communication mechanisms at each
level of the hierarchy.

\textbf{Level $0 \to 1$ (cell $\to$ tissue coupling).}
This is the primary coupling layer and the level at which most
cancers sever. The coupling mechanisms are:
\begin{itemize}
\item \textbf{Gap junctions} (connexins, especially Cx43): direct
   cytoplasmic channels between adjacent cells that share ions,
   metabolites, and second messengers. Gap junctional intercellular
   communication (GJIC) is the highest-bandwidth cell-tissue coupling
   channel. Connexin-43 downregulation is among the earliest events
   in many epithelial cancers \cite{Aasen2016}.
\item \textbf{E-cadherin} (CDH1): homophilic adhesion molecule that
   mediates contact inhibition of proliferation. Loss of E-cadherin
   is the defining event of epithelial-mesenchymal transition (EMT)
   and directly severs the contact-dependent component of $\lambda_1$
   \cite{Hanahan2011}.
\item \textbf{Juxtacrine signaling} (Notch--Delta): direct
   cell-cell contact signaling that maintains tissue identity and
   differentiation state. Requires physical adjacency, destroyed
   by loss of tissue architecture.
\item \textbf{Integrin--ECM attachment}: anchorage to the
   extracellular matrix and basement membrane. Loss of anchorage
   dependence (anoikis resistance) severs the structural component
   of tissue coupling.
\end{itemize}

\textbf{Level $1 \to 2$ (tissue $\to$ organ coupling).}
\begin{itemize}
\item \textbf{Paracrine signaling}: local growth factors and
   morphogen gradients (Wnt, Hedgehog, BMP, TGF-$\beta$) that
   coordinate tissue-level organization within the organ context.
\item \textbf{Extracellular matrix architecture}: the structural
   scaffold that encodes tissue-level spatial organization and
   transmits mechanical signals.
\item \textbf{Local vasculature}: nutrient and oxygen delivery
   that ties tissue metabolism to organ-level regulation. Vascular
   normalization restores this coupling channel.
\end{itemize}

\textbf{Level $2 \to 3$ (organ $\to$ organism coupling).}
\begin{itemize}
\item \textbf{Endocrine signaling}: systemic hormones (insulin,
   estrogen, cortisol, thyroid hormone) that couple organ function
   to organismal state. Hormone-dependent cancers (breast, prostate)
   retain partial $\lambda_3$ coupling, explaining their relative
   treatability.
\item \textbf{Immune surveillance}: T-cells, NK cells, and
   macrophages that enforce organism-level quality control on
   aberrant cells. Immune checkpoint molecules (PD-L1, CTLA-4) are
   the molecular implementation of this coupling channel.
\item \textbf{Neural innervation}: autonomic nervous system
   regulation of organ function and tissue homeostasis.
\item \textbf{Circadian regulation}: organism-level temporal
   coordination via the SCN-driven hormonal cascade.
\end{itemize}

More precisely, each aggregate $\lambda_j$ decomposes into per-channel
components. Chain severance can be partial: a cell may lose
gap junctional coupling ($\lambda_1^{\mathrm{GJIC}} \to 0$) while
retaining cadherin-mediated contact inhibition
($\lambda_1^{\mathrm{CDH1}} > 0$). This partial severance produces
the intermediate phenotypes observed in pre-malignant lesions.
Complete severance across all channels at a given level produces the
full malignant phenotype.

% ------------------------------------------------------------------
\section{Main Result}
\label{sec:main}
% ------------------------------------------------------------------

\begin{proposition}[Escalation Chain Severance Produces Convergent
Malignancy]
\label{thm:cancer}
Let $E_0$ be a cell in a multicellular organism with intact
escalation chain $E_0 \to E_1 \to \cdots \to E_m$, modeled as an
entity under the bounded-context framework. If chain severance
occurs at level $k = 0$ (cell loses coupling to tissue), then:

\textbf{Part I.} $E_0$ optimizes under $C_n^{(0)}$ alone.
By graph necessity \cite{McCarthy2026a}, it maintains only the
propositions serviceable within cell-level bounded context.

\textbf{Part II.} By $K/A$ inseparability \cite{McCarthy2026c}, the
knowledge base $K_0$ contracts to index only cell-level actions
$A_0$: those executable without tissue-level coordination.

\textbf{Part III.} By the escalation convergence result \cite{McCarthy2026b}, with
no higher level to escalate to, all optimization pressure concentrates
at $E_0$. The cell becomes the terminal level of its own hierarchy.

\textbf{Part IV.} The action space under cell-level $C_n^{(0)}$
converges across cell types. $C_n^{(0)}$ is small relative to the
full repertoire of cell-level actions, so the asymmetric retention
criterion of \cite{McCarthy2026c} forces displacement toward the
same survival-essential set: with limited slots and the same
survival-critical actions competing for them, tissue-specific content
is progressively displaced by actions with nonzero expected benefit
under cell-level optimization.
\end{proposition}

\textbf{Argument.}

\textbf{Part I: Context Collapse.}

Before severance, $E_0$ operates with effective context
$C_n^{\mathrm{eff}} = C_n^{(0)} + \sum_{j=1}^{m} \lambda_j
C_n^{(j)}$. The cell maintains propositions that reference
tissue-level and organism-level states: differentiation signals,
growth factor dependencies, contact inhibition cues, apoptotic
triggers from the tissue microenvironment. These propositions are
serviceable \emph{because} the escalation chain provides the context
in which they have grounding.

After severance ($\lambda_j = 0$ for $j \geq 1$), these
propositions lose their grounding: they reference states that $E_0$
can no longer observe (tissue integrity, organ function, organismal
need). By graph necessity \cite{McCarthy2026a}, propositions without
grounding in observable states cannot be maintained under bounded
context without displacing grounded propositions. The retention
criterion of \cite{McCarthy2026c} -- retain when benefit is
uncertain, prune when benefit is certainly zero -- applies: the
benefit of tissue-referencing propositions is certainly zero to a
cell that cannot observe tissue state. They are pruned.

$K_0$ contracts to propositions grounded in cell-level observables:
local nutrient concentration, immediate mechanical environment,
intracellular state.

\textbf{Part II: Action Space Contraction.}

By $K/A$ inseparability \cite{McCarthy2026c}, $K_0$ indexes $A_0$
and vice versa. As $K_0$ contracts to cell-level propositions, $A_0$
contracts to cell-level actions. Actions that require tissue-level
coordination -- responding to contact inhibition, executing
apoptosis on organismal command, maintaining differentiated quiescence
-- lose their knowledge support. They are not ``disabled'' by
mutation; they become \emph{inaccessible} because the knowledge
required to execute them (tissue state, organismal need) is no longer
in $K_0$.

What remains in $A_0$ are actions executable with cell-level knowledge
alone:
\begin{itemize}
\item \textbf{Proliferation}: cell division requires only local
   nutrient availability and intracellular state; no tissue-level
   coordination needed.
\item \textbf{Metabolic autonomy}: aerobic glycolysis (the Warburg
   effect \cite{Warburg1956}) is the metabolic mode with the fewest
   tissue-level \emph{dependencies}: it does not require reliable
   tissue-level oxygen delivery infrastructure. Critically, cancer
   cells use glycolysis even when oxygen is present---not because
   oxygen is unavailable, but because after severance the cell lacks
   tissue-level information about oxygen supply reliability and
   defaults to the metabolic strategy that requires no such
   information. This does not imply mitochondrial shutdown (many
   cancer cells retain functional oxidative phosphorylation and use
   both pathways), but glycolytic capacity provides metabolic
   independence from tissue-level coordination.
\item \textbf{Motility}: movement in response to local nutrient
   gradients requires only cell-level sensing.
\item \textbf{Immune evasion}: both passive and active mechanisms.
   MHC-I surface markers, maintained by tissue-level signaling, decay
   when that signaling is lost. Additionally, cancer cells actively
   upregulate PD-L1 in response to immune attack (IFN-$\gamma$
   signaling)---an adaptive response that requires only cell-level
   sensing of immune pressure, not tissue-level coordination.
\end{itemize}

\textbf{Part III: Terminal Optimization.}

In the intact organism, the cell's optimization is bounded by
escalation: problems exceeding $C_n^{(0)}$ escalate to tissue, organ,
or organism level. The cell does not need to solve every problem;
it only needs to handle its local scope and signal failure upward.

After severance, there is no higher level. By the escalation
convergence result \cite{McCarthy2026b}, this means all optimization
pressure concentrates at $E_0$. The cell is no longer a component
in a hierarchy; it is a complete autonomous entity. Its optimization
target is its own persistence, gradient descent on $U$ at cell
level \cite{McCarthy2026d}, with no constraint from above.

This is precisely the optimization regime of a unicellular organism:
an autonomous entity with cell-level $C_n$, optimizing for its own
survival under bounded context. The cell has not ``reverted'' to a
unicellular program. It has been \emph{forced into the same
optimization regime} by the loss of its escalation chain.

\textbf{Part IV: Convergence.}

The convergence argument is not that all severed cells must have
identical knowledge graphs. Graph necessity \cite{McCarthy2026a}
requires graph \emph{structure}, not identical graph \emph{content}.
The convergence is at the level of the action space, and it follows
from the interaction of bounded context with the asymmetric retention
criterion.

$C_n^{(0)}$ is small relative to the full repertoire of possible
cell-level actions. Under the displacement criterion of
\cite{McCarthy2026c}, when $K_0$ is at capacity, a proposition $p$
is displaced by $q$ only when $\mathbb{E}[G(a_q, K)] >
\mathbb{E}[G(a_p, K)]$, strictly higher marginal survival value.
After severance, tissue-specific propositions have zero marginal
value (they reference unobservable states), while cell-level
survival actions (proliferation, nutrient acquisition, damage repair)
have nonzero value regardless of tissue of origin. The retention
criterion therefore drives displacement in the same direction for
all severed cells: tissue-specific content is replaced by cell-level
survival content.

The convergence is not perfect; residual tissue-specific $K_0$
content produces tissue-specific variation, especially early in
progression. But the attractor is the same: the set of actions with
highest marginal survival value under $C_n^{(0)}$. This set is
constrained by cell biology, not by tissue of origin, so it
converges across cancer types.

This offers a structural account for why Hanahan and Weinberg's
hallmarks \cite{Hanahan2011} are observed across cancer types: they
are not shared mutations or shared ancient programs, but the
\emph{convergent action space} that the retention criterion selects
under cell-level bounded context.
\medskip

\textbf{Status of this result.} Proposition~\ref{thm:cancer} is a
conceptual argument, not a formal derivation. Each Part invokes a
result from the formal chain (Papers~1--4) and applies it to the
biological setting by analogy. The analogies are specific and
testable---Section~\ref{sec:predictions} states falsifiable
predictions---but the Proposition does not prove that the hallmarks
are the \emph{unique} convergent action set under cell-level $C_n$.
It argues that convergence to a small, survival-essential set is
structurally predicted, and that the observed hallmarks are consistent
with that prediction.

% ------------------------------------------------------------------
\section{Corollaries}
\label{sec:corollaries}
% ------------------------------------------------------------------

\begin{corollary}[Velocity of Malignant Transformation]
\label{cor:velocity}
Malignant transformation does not require the accumulation of a
specific set of mutations. It requires chain severance: the
loss of coupling to the tissue-level escalation chain. Chain
severance can occur through a single event (loss of a key signaling
receptor, disruption of gap junctions, escape from the tissue
microenvironment) rather than through the sequential accumulation
of independent mutations.
\end{corollary}

\textbf{Argument.}
The somatic mutation theory requires multiple independent hits
\cite{Vogelstein2013}: each hallmark capability must be separately
acquired through mutation. Under our framework, the hallmarks are
not independently acquired; they are jointly \emph{released}
when the escalation chain severs. A single event that decouples the
cell from tissue-level signaling simultaneously enables all
hallmark behaviors, because all of them are cell-level actions that
were \emph{suppressed} by tissue-level coordination, not
\emph{absent} from the cell's intrinsic repertoire.

This explains why cancer can progress rapidly from a single
initiating event: the ``multi-hit'' appearance is a consequence of
progressive chain severance (each hit weakens coupling further)
rather than independent capability acquisition. The velocity is
structural, not stochastic.

\textbf{Compatibility with multi-hit observations.} The framework
does not contradict the observation that multiple mutations are
typically found in cancers. Rather, it reinterprets their role:
mutations are mechanisms of chain severance (disrupting signaling
pathways, disabling apoptotic responses to tissue-level commands),
not independent capability acquisitions. The ``multi-hit'' pattern
reflects progressive decoupling: each mutation weakens a different
coupling channel $\lambda_j$, rather than independent construction
of each hallmark.
\medskip

\begin{corollary}[Therapeutic Resistance]
\label{cor:resistance}
Established tumors exhibit persistent therapeutic resistance because
the severed cell is \emph{locally rational}: it is optimizing
correctly under its own bounded context.
\end{corollary}

\textbf{Argument.}
Therapies that target specific mutations (targeted therapy) or
specific proliferative mechanisms (chemotherapy) impose selection
pressure on an entity that is already performing gradient descent on
$U$ at cell level \cite{McCarthy2026d}. The cell's $\gamma > 0$
(Paper~4, \cite{McCarthy2026d}) guarantees continued engagement with
adaptive mechanisms. Its learning efficiency $\eta_M > 0$
(the rate of improvement per unit of learning signal;
\cite{McCarthy2026d}) guarantees learning from
therapeutic pressure. The therapy becomes a new environmental signal
that the cell incorporates into its optimization.

This explains why resistance is the norm rather than the exception:
the cell is not a broken machine to be fixed but an autonomous
entity optimizing under its own constraints. Therapy that does not
restore the escalation chain merely changes the optimization
landscape without changing the optimization regime.

Effective therapy, on this account, must either:
\begin{enumerate}
\item \textbf{Restore the chain}: re-establish tissue-level coupling
   so that $C_n^{\mathrm{eff}} > C_n^{(0)}$, returning tissue-level
   actions to the cell's accessible action space. Differentiation
   therapy (e.g., all-trans retinoic acid in acute promyelocytic
   leukemia) operates this way.
\item \textbf{Exploit the constraint}: target cell-level actions
   that are \emph{necessary} under $C_n^{(0)}$, not just contingently
   adopted. If an action is necessary under cell-level bounded
   context, the cell cannot evolve away from it without leaving the
   $C_n^{(0)}$ optimization regime.
\end{enumerate}
\medskip

\begin{corollary}[The Atavistic Pattern Without Atavistic Causation]
\label{cor:atavistic}
The phylostratigraphic signal -- cancer genes being
disproportionately ancient \cite{Trigos2017} -- is real but does
not imply atavistic reversion. It reflects convergent necessity.
\end{corollary}

\textbf{Argument.}
Genes that are ancient are precisely the genes that were optimized
for cell-level bounded context during the unicellular era. When a
modern cell undergoes chain severance and is forced back to
cell-level optimization, the \emph{same genes} become relevant
again, not because an ancient program is ``reactivated'' but
because the same constraint ($C_n^{(0)}$) selects for the same
solutions.

The atavistic theory \cite{Davies2011} correctly identifies the
phenotypic pattern (cancer cells resemble unicellular organisms) and
the genetic evidence (ancient genes are upregulated). Its error is
causal: it infers historical reversion where the mechanism is
structural convergence. The genes are not activated \emph{because}
they are ancient; they are activated because they encode the actions
that are necessary under cell-level bounded context, and they happen
to be ancient because unicellular organisms faced the same constraint.

This distinction is testable: the atavistic theory predicts that the
ancient gene set forms a coherent ``unicellular program'' that is
activated as a unit. Our framework predicts that the set of genes
activated in any given cancer will be determined by the specific
$C_n^{(0)}$-level optimization landscape, and will include ancient
genes \emph{to the extent that} they encode cell-level actions,
but will also include recent genes that happen to serve cell-level
optimization (e.g., genes co-opted for immune evasion that evolved
in the multicellular era).
\medskip

\begin{corollary}[Metastasis as Optimal Foraging Under Cell-Level
$C_n$]
\label{cor:metastasis}
Metastasis is not an acquired ``capability'' but the default
exploratory behavior of an autonomous entity optimizing under
bounded context in a spatially heterogeneous environment.
\end{corollary}

\textbf{Argument.}
In an intact organism, cell motility is constrained by tissue-level
signaling: basement membrane attachment, contact inhibition, tissue
boundary maintenance. These constraints are maintained by the
escalation chain. The tissue level enforces spatial organization
that individual cells could not compute independently.

After chain severance, the cell has no spatial constraint beyond its
local environment. By the non-termination result of
\cite{McCarthy2026d} (Theorem~1, Part~II), the reachable trajectory
space $T_{\mathrm{reachable}}(K_0)$ grows monotonically with $K_0$.
A cell in a nutrient-poor local environment will, through gradient
descent on $U$, discover that movement to nutrient-rich regions
reduces uncertainty about survival.

\textbf{Inefficiency is predicted, not anomalous.} Metastasis is
exploration under $C_n^{(0)}$---cell-level bounded context with no
tissue-level guidance. More than 99.9\% of circulating tumor cells
die without colonizing a distant site \cite{Massague2016}. This
massive inefficiency is exactly what the framework predicts:
exploration without higher-level coordination is high-variance.
The cell has $\gamma > 0$ (it engages adaptive mechanisms) but
its context window is too small to model the circulatory environment,
immune surveillance, or distant-site compatibility. Each step is
locally rational---movement reduces local uncertainty---but the
cell lacks the $C_n^{\mathrm{eff}}$ to plan a viable multi-step
trajectory. The rare successful metastases are not evidence of
acquired capability but survivorship bias over massive parallel
exploration.

The ``invasion-metastasis cascade'' of epithelial-mesenchymal
transition, intravasation, survival in circulation, and colonization
at distant sites is not a sequence of independently acquired
capabilities. It is the progressive unfolding of cell-level
exploration in a spatially structured environment, where each step
reveals new states that reduce local uncertainty.
\medskip

\begin{corollary}[Transplant Rejection as Chain Mismatch]
\label{cor:transplant}
Organ transplant rejection is not chain severance but \emph{chain
mismatch}: the transplanted tissue has intact internal coupling
(levels $0 \to 1$) but mismatched organism-level coupling
(level $2 \to 3$). The immune system's rejection response is the
escalation chain functioning correctly, detecting cells whose
coupling markers do not match the recipient organism.
\end{corollary}

\textbf{Argument.}
A transplanted organ retains its internal tissue organization: the
cells are coupled to each other and to tissue-level architecture
via gap junctions, cadherins, and paracrine signaling. The
$\lambda_1$ coupling channels are intact. What is mismatched is
$\lambda_3^{\mathrm{immune}}$: the organism-level immune
surveillance system recognizes the transplanted cells' MHC markers
as foreign, not coupled to \emph{this} organism's escalation
chain.

At the level of the coupling variable $\lambda_3^{\mathrm{immune}}$,
chain severance (cancer) and chain mismatch (transplant) produce the
same signal: absent or aberrant organism-level coupling. The immune
mechanisms differ in detail---transplant rejection is driven by MHC
mismatch recognition, while cancer immune evasion involves PD-L1
upregulation and MHC-I downregulation---but structurally, both
represent failures of the level $2 \to 3$ coupling, and the
organism-level enforcement mechanism responds to both.

This reframes the immunosuppression tradeoff structurally.
Immunosuppressive therapy (cyclosporine, tacrolimus) weakens the
organism-level coupling enforcement
$\lambda_3^{\mathrm{immune}} \to 0$ globally. This permits the
mismatched organ but simultaneously disables the surveillance
mechanism that detects genuine chain severance. The well-documented
2--5$\times$ elevation in cancer incidence among transplant
recipients \cite{Engels2011} is the predicted consequence:
weakened chain enforcement permits both the intended mismatch
\emph{and} unintended severance events.

\textbf{Tolerance induction} -- protocols that establish mixed
hematopoietic chimerism or regulatory T-cell-mediated acceptance
-- operates differently. Rather than weakening enforcement, it
re-establishes matching coupling: the recipient's immune system
learns to recognize the transplanted tissue as part of its own
escalation chain. The coupling
$\lambda_3^{\mathrm{immune}}$ is restored for the transplant
specifically, while global enforcement remains intact. The
framework predicts that tolerance induction should produce better
long-term outcomes than chronic immunosuppression, not only
for graft survival but for overall survival, because it
preserves the organism-level chain enforcement that prevents
cancer. Early clinical evidence from mixed chimerism protocols
is consistent with this prediction \cite{Kawai2008}.

HLA matching reduces the mismatch at the coupling interface itself,
minimizing the signal that triggers enforcement without weakening
enforcement. This is why HLA-matched transplants require less
immunosuppression and have better outcomes: less mismatch means
less enforcement to suppress.

\textbf{Coupling-channel specificity predicts cancer type distribution.}
Immunosuppression weakens $\lambda_3^{\mathrm{immune}}$ specifically
---organism-level immune enforcement---while leaving tissue-level
coupling ($\lambda_1$: gap junctions, contact inhibition, basement
membrane integrity) intact. The framework therefore predicts that
immunosuppressed patients should show disproportionately elevated
rates of cancers where immune surveillance is the \emph{primary}
coupling channel, and near-baseline rates of cancers where tissue
architecture provides the dominant restraint.

This is confirmed by the transplant cancer data \cite{Engels2011}.
Cancers dominated by immune-surveillance coupling show massive
elevation: Kaposi sarcoma (${\sim}60\times$), non-Hodgkin lymphoma
(${\sim}7\times$), and non-melanoma skin cancers---all
immune-surveilled or virally mediated. Cancers where tissue
architecture is the primary coupling mechanism show minimal
elevation: breast (${\sim}1.2\times$), prostate (${\sim}0.9\times$),
and colorectal (${\sim}1.2\times$)---cancers restrained primarily
by glandular organization, basement membrane, and epithelial
architecture, none of which immunosuppression disrupts.

This pattern is difficult to explain under the somatic mutation
theory (which predicts uniform elevation proportional to reduced
immune clearance of mutant cells) or the atavistic theory (which
has no mechanism for coupling-channel specificity). It follows
directly from the framework: weakening one coupling channel
($\lambda_3^{\mathrm{immune}}$) permits severance only in tissues
where that channel was load-bearing.
\medskip

% ------------------------------------------------------------------
\section{Mechanisms of Chain Severance}
\label{sec:mechanisms}
% ------------------------------------------------------------------

Chain severance is not a single event but a spectrum of decoupling
mechanisms, each reducing $\lambda_j$ toward zero:

\textbf{Genetic mechanisms.} Mutations in signaling receptors (e.g.,
E-cadherin loss), tumor suppressor pathways (e.g., p53, Rb), or
growth factor dependencies reduce the cell's capacity to receive
tissue-level signals. These are the mechanisms identified by SMT
\cite{Vogelstein2013}. In our framework, they are not causes of
cancer but \emph{mechanisms of chain severance}: each mutation
weakens a specific coupling channel.

\textbf{Microenvironmental mechanisms.} Chronic inflammation,
fibrosis, or tissue damage disrupts the signaling infrastructure
through which higher levels communicate with cells. This is the
mechanism identified by TOFT \cite{Sonnenschein2000}: the tissue
architecture that maintains coupling is destroyed, and cells
decouple even without intrinsic mutations. The coupling
$\lambda_j \to 0$ occurs not because the cell's receivers are
broken but because the signal is no longer transmitted.

\textbf{Epigenetic mechanisms.} Methylation changes, chromatin
remodeling, or transcriptional reprogramming can silence the
genes encoding tissue-level signaling components without genomic
mutation. This produces chain severance through altered gene
expression rather than altered sequence.

\textbf{Spatial mechanisms.} Physical separation from the tissue
microenvironment -- through basement membrane breach, for
instance -- removes the cell from the signaling field. The cell's
coupling hardware is intact; the coupling medium is absent.

All four mechanisms produce the same outcome: $\lambda_j \to 0$,
effective context collapse to $C_n^{(0)}$, and optimization under
cell-level bounded context. The convergent phenotype (the hallmarks)
follows regardless of \emph{which} mechanism severed the chain.
This explains why cancers arising from different mechanisms
(hereditary, environmental, sporadic) converge on the same
phenotypic endpoints.

% ------------------------------------------------------------------
\section{The Encoding Hierarchy in Cancer}
\label{sec:encoding}
% ------------------------------------------------------------------

The encoding hierarchy of \cite{McCarthy2026d} has a specific
consequence for cancer that connects to the empirical phenomenon of
tumor heterogeneity.

In the intact organism, actions at each encoding level have been
validated to the threshold $\theta_i = 1 - \varepsilon / C_i$. The
genomic level ($L_0$) encodes actions validated across evolutionary
timescales. The epigenetic level ($L_1$) encodes tissue-specific
programs validated across developmental timescales. The active
context level ($L_n$) encodes moment-to-moment responses.

After chain severance, the encoding hierarchy is disrupted at all
levels above the cell:

\begin{itemize}
\item \textbf{Organism-level encodings} (hormonal responses,
   circadian regulation, immune cooperation) lose their validation
   context. They were validated under $C_n^{\mathrm{eff}}$; they are
   now evaluated under $C_n^{(0)}$. Many fail the retention criterion.

\item \textbf{Tissue-level encodings} (differentiation state,
   contact inhibition, growth factor dependence) similarly lose
   validation context.

\item \textbf{Cell-level encodings} (basic metabolism, division
   machinery, DNA repair) retain their validation. They were
   validated under $C_n^{(0)}$ originally and remain valid.
\end{itemize}

The result is \emph{selective devalidation}: the encoding hierarchy
does not collapse uniformly but selectively loses its higher-level
content. What remains is a cell operating on cell-level encodings,
the same encodings that a unicellular organism would maintain. This
is the encoding-level explanation for the atavistic gene expression
pattern observed by Trigos et al.\ \cite{Trigos2017}: ancient genes
are preferentially retained not because they are ``reactivated'' but
because they are the ones whose validation survives context collapse.

\textbf{Tumor heterogeneity.} The progressive nature of context
collapse explains intratumor heterogeneity. Cells at different
stages of chain severance retain different amounts of higher-level
encoding. Early-stage tumor cells may retain partial tissue-level
coupling (some $\lambda_j > 0$), producing a mixed phenotype. As
severance progresses, the phenotype converges toward the cell-level
optimum. This gradient of severance within a single tumor produces
the phenotypic heterogeneity observed empirically, without requiring
different clonal origins for different phenotypes.

Clonal evolution \cite{Nowell1976} describes the \emph{dynamics}
of this process: cells with more complete severance have higher
fitness in the decoupled regime and are selected. But the direction
of evolution is not random; it converges on the cell-level
optimum, which is structurally determined by $C_n^{(0)}$.

% ------------------------------------------------------------------
\section{Engagement with Prior Work}
\label{sec:prior}
% ------------------------------------------------------------------

\textbf{Hanahan and Weinberg \cite{Hanahan2000, Hanahan2011}:
Hallmarks of Cancer.}
The hallmarks -- sustaining proliferative signaling, evading growth
suppressors, resisting cell death, enabling replicative immortality,
inducing angiogenesis, activating invasion and metastasis, plus the
enabling characteristics and emerging hallmarks of the 2011 update
-- are described as ``capabilities'' that must be ``acquired.'' Our
framework reinterprets them: they are not acquired but
\emph{revealed}. They are the necessary action space of any entity
optimizing under cell-level bounded context. The hallmarks are
universal across cancer types because $C_n^{(0)}$ is universal
across cell types. What differs across cancers is the \emph{route}
of chain severance (which coupling channels fail first), not the
\emph{destination} (cell-level optimization).

\textbf{Davies and Lineweaver \cite{Davies2011}: Atavistic Theory.}
See Corollary~\ref{cor:atavistic}. The atavistic theory correctly
identifies the phenotypic and phylostratigraphic pattern. Our
framework provides the mechanism it lacks (chain severance) and
corrects the causal inference (convergent necessity, not historical
reversion). The two theories make different predictions about which
genes will be activated in any given cancer: the atavistic theory
predicts a coherent ancient program; our framework predicts
gene activation determined by $C_n^{(0)}$-level optimization,
which overlaps heavily but not completely with the ancient gene set.

\textbf{Sonnenschein and Soto \cite{Sonnenschein2000}: TOFT.}
TOFT is the closest existing framework to our account. It correctly
identifies the organizational level (tissue, not cell) at which
cancer originates, and correctly argues that proliferation is the
default state of cells, with tissue organization as the restraint.
Our framework provides TOFT with the formalism it has lacked: the
escalation chain is the formal structure of tissue organization,
$C_n^{\mathrm{eff}}$ is the formal measure of organizational
integration, and chain severance is the formal description of
organizational breakdown. TOFT's ``default proliferative state'' is
our ``necessary action space under cell-level $C_n$.''

\textbf{Nowell \cite{Nowell1976}: Clonal Evolution.}
Clonal evolution describes the dynamics of cancer progression as
Darwinian selection among cell populations. Our framework does not
replace this account but provides its \emph{attractor}: clonal
evolution does not wander randomly through phenotype space but
converges on the cell-level optimum determined by $C_n^{(0)}$. The
direction of clonal evolution is structurally determined; the
stochastic component is in \emph{which cells} reach the attractor
first, not \emph{where} the attractor is.

\textbf{Huang et al.\ \cite{Huang2009}: Attractor Landscape.}
Huang's gene-regulatory-network framework models cell states as
attractors in a high-dimensional expression landscape, with cancer
as occupation of abnormal attractors that normal cells rarely visit.
Our framework subsumes this picture: the attractors that Huang
identifies are the \emph{stable configurations} of cell-level
optimization under $C_n^{(0)}$. Chain severance does not create new
attractors; it removes the tissue-level coupling that steered cells
away from the cell-level optimum. What Huang describes as an
``abnormal attractor'' is, on our account, the \emph{default}
attractor under cell-level bounded context---abnormal only relative
to tissue-integrated operation. The two frameworks make compatible
predictions, but ours provides a principled reason \emph{why} the
attractor landscape has the structure it does: the bounded-context
constraint $C_n^{(0)}$ determines which configurations are stable
for an autonomous cell.

\textbf{Aktipis et al.\ \cite{Aktipis2015}: Cancer as Cheating.}
The evolutionary cooperation framework describes cancer as a
breakdown of multicellular cooperation: cells that ``cheat'' by
proliferating at the expense of the collective. Our framework
provides the structural basis: cooperation is maintained by the
escalation chain (higher organizational levels enforce cooperative
behavior by extending $C_n^{\mathrm{eff}}$). ``Cheating'' is not a
strategy adopted by the cell but the necessary outcome of chain
severance: a cell that cannot observe tissue-level state cannot
cooperate with the tissue, because cooperation requires the
knowledge base ($K_0$) to contain tissue-referencing propositions
that are no longer grounded.

\textbf{Merlo et al.\ \cite{Merlo2006}: Cancer as Ecological Process.}
The ecological framework describes tumor progression as an
evolutionary process within the ecosystem of the body. Our framework
adds the structural constraint: the ``ecosystem'' is the escalation
chain, and cancer progression is the progressive collapse of that
chain. The ecological metaphor is correct but incomplete; it
describes dynamics without explaining why those dynamics converge.

% ------------------------------------------------------------------
\section{The Central Prediction: Coupling-Channel Specificity}
\label{sec:predictions}
% ------------------------------------------------------------------

The framework's unique predictive signature is \textbf{coupling-channel
specificity}: the cancer type that emerges from any disruption is
determined by \emph{which coupling channel} $\lambda_j$ was
load-bearing for that tissue. This is the one prediction that
distinguishes the chain-severance framework from all existing cancer
theories. The somatic mutation theory predicts cancer incidence based on
proliferation rate and mutational exposure. The atavistic theory
predicts a uniform reversion to unicellular programs. Neither has a
mechanism for predicting which cancer type emerges from which
disruption. The chain-severance framework does: sever a specific
channel, and the cancers that appear are those where that channel was
the primary organizational restraint.

\subsection{Channel-specific natural experiments}

Five classes of natural experiment test this principle, each disrupting
a different coupling channel and producing a different cancer type
distribution (Table~\ref{tab:channels}).

\begin{table}[ht]
\centering
\small
\begin{tabular}{p{2.8cm} p{3.2cm} p{3.8cm} p{2.8cm}}
\hline
\textbf{Disruption} & \textbf{Channel weakened} & \textbf{Predicted cancers} & \textbf{Evidence} \\
\hline
Immunosuppression (transplant) &
$\lambda_3^{\mathrm{immune}}$: organism-level immune surveillance &
Immune-surveilled: lymphoma (${\sim}7\times$), Kaposi (${\sim}60\times$), skin.
Not tissue-architecture cancers (breast ${\sim}1.2\times$, prostate ${\sim}0.9\times$). &
\cite{Engels2011} \\[6pt]
Foreign body implant &
$\lambda_1$: local tissue architecture (basement membrane, contact inhibition) &
Mesenchymal (sarcomas) at implant site. Not lymphomas. &
\cite{Brand1975} \\[6pt]
Denervation &
$\lambda_2^{\mathrm{neural}}$: neural coupling &
Cancers in denervated tissue where neural signaling maintains differentiation. &
\cite{Zhao2017} \\[6pt]
Chronic inflammation (Barrett's, hepatitis, IBD) &
$\lambda_1 + \lambda_3^{\mathrm{immune}}$: tissue architecture and immune signaling &
Site-specific cancers at inflammation locus: esophageal, hepatocellular, colorectal. &
Well established \\[6pt]
Endocrine disruption &
$\lambda_3^{\mathrm{endocrine}}$: hormonal coupling &
Hormone-responsive cancers (breast, prostate, endometrial). Not cancers where hormonal signaling is not load-bearing. &
\cite{Herbst1971} \\
\hline
\end{tabular}
\caption{Coupling-channel specificity: disrupting different
$\lambda_j$ channels produces different cancer type distributions.
Each row is a natural experiment testing the framework's central
prediction.}
\label{tab:channels}
\end{table}

\subsection{Carcinogen classification by coupling channel}

The principle extends to carcinogenesis. Two carcinogens hitting the
same tissue but damaging different coupling channels should produce
different cancer subtypes. Lung cancer provides the clearest test.
Smoking causes chronic inflammation and epithelial remodeling---damage
to $\lambda_1$ (tissue architecture)---and is associated primarily
with squamous cell carcinoma and adenocarcinoma. Radon causes direct
cell-level DNA damage via alpha particles and is associated primarily
with small cell lung cancer: poorly differentiated, neuroendocrine,
maximally aggressive. Different channel, different subtype, same
organ.

The prediction generalizes: \textbf{carcinogen classification should
specify which coupling channel the carcinogen disrupts, not merely
which tissue is exposed, because the disrupted channel determines
the cancer subtype.} This is testable by re-analyzing existing
epidemiological datasets that record both exposure type and cancer
subtype within the same organ.

\subsection{Comorbidity as coupling-channel degradation}

The coupling-channel principle extends beyond exogenous carcinogens
to endogenous disease. Any disease that chronically degrades a
specific coupling channel~$\lambda_j$ should increase the rate of
cancers where $\lambda_j$ is load-bearing---and should increase the
\emph{specific subtype} associated with that channel, not merely
cancer incidence in general. Table~\ref{tab:comorbidity} compiles
six well-established comorbidity associations, each currently treated
as a separate finding in a separate literature. The framework unifies
them under a single principle.

\begin{table}[ht]
\centering
\small
\caption{Comorbidity-cancer associations stratified by
coupling-channel degradation. Each disease degrades a specific
coupling channel; the cancers that increase are those where that
channel is load-bearing. All associations are well-established
epidemiologically but are not currently unified under a common
mechanism.}
\label{tab:comorbidity}
\begin{tabular}{@{}p{3.0cm}p{3.0cm}p{3.5cm}p{3.0cm}@{}}
\toprule
\textbf{Disease} & \textbf{Channel degraded} &
\textbf{Cancer increase} & \textbf{Channel match} \\
\midrule
HIV/AIDS &
$\lambda_3^{\mathrm{immune}}$ (immune surveillance) &
Kaposi's, NHL, cervical &
Immune-coupled \\[4pt]
Diabetes / metabolic syndrome &
$\lambda_2^{\mathrm{metabolic}}$ (metabolic coupling) &
Pancreatic, liver, endometrial &
Metabolic-coupled \\[4pt]
Vagotomy &
$\lambda_2^{\mathrm{neural}}$ (neural signaling) &
Gastric &
Neural-coupled \\[4pt]
Inflammatory bowel disease &
$\lambda_1$ (tissue architecture) &
Colorectal &
Architecture-coupled \\[4pt]
Cirrhosis &
$\lambda_1$ (tissue architecture) &
Hepatocellular carcinoma &
Architecture-coupled \\[4pt]
PCOS / endocrine disorders &
$\lambda_3^{\mathrm{endocrine}}$ (hormonal coupling) &
Endometrial, breast &
Endocrine-coupled \\
\bottomrule
\end{tabular}
\end{table}

The diagnostic prediction is immediate: \textbf{knowing which
coupling channel a disease degrades predicts which cancer subtypes
will increase in that patient population.} This is testable by
cross-referencing disease-specific cohorts (e.g., diabetes registries,
autoimmune disease registries) with cancer subtype incidence, stratified
by coupling-channel assignment. Existing comorbidity databases are
sufficient; the framework provides the organizing variable that is
currently missing.

The converse prediction also holds: diseases that \emph{strengthen}
a coupling channel should be protective against the corresponding
cancer subtypes. Autoimmune conditions---characterized by
\emph{overactive} immune coupling---should protect against cancers
where immune surveillance is the load-bearing channel, while
potentially increasing risk for cancers driven by chronic
inflammatory tissue damage. This bidirectional prediction
distinguishes the framework from generic ``inflammation causes
cancer'' accounts, which predict only increased risk.

\subsection{Specific predictions}

The coupling-channel principle generates the following specific,
falsifiable predictions:

\begin{enumerate}
\item \textbf{Coupling measurement predicts malignancy.} Quantifying
   $\lambda_j$ should predict malignant potential more accurately than
   mutational burden alone. Cells with low $\lambda_j$ but few
   mutations should progress to malignancy; cells with high
   mutational burden but intact $\lambda_j$ should not. Testable by
   measuring intercellular communication (gap junction activity,
   paracrine signaling responsiveness, contact inhibition) alongside
   mutational profiling.

\item \textbf{Convergent gene expression is constraint-determined,
   not phylogeny-determined.} Cancer gene expression should converge
   on a set determined by cell-level optimization under $C_n^{(0)}$,
   which \emph{mostly but not entirely} overlaps with
   phylogenetically ancient genes. Recently evolved genes that serve
   cell-level autonomous function (e.g., immune evasion mechanisms
   from the multicellular era) should be upregulated despite being
   phylogenetically recent. The atavistic theory predicts they
   should not be.

\item \textbf{Differentiation therapy should outperform cytotoxic
   therapy.} Therapies that restore tissue-level coupling should
   produce more durable responses than therapies that kill cells
   without restoring coupling, because the former addresses the
   structural cause (channel severance) while the latter imposes
   selection pressure on a locally rational optimizer.

\item \textbf{Tumor heterogeneity correlates with severance
   gradient.} Within a single tumor, cells with residual
   tissue coupling ($\lambda_j > 0$) should exhibit more
   differentiated phenotypes than fully decoupled cells. Testable by
   spatial transcriptomics combined with intercellular communication
   assays.

\item \textbf{Necessary hallmarks vs.\ contingent hallmarks.}
   Some hallmarks are \emph{necessary} under cell-level $C_n$
   (proliferation, metabolic autonomy); others are \emph{contingent}
   on the specific optimization landscape (angiogenesis induction,
   specific immune evasion strategies). Necessary hallmarks should
   appear in every cancer; contingent hallmarks should vary by
   context and by which coupling channel was severed.

\item \textbf{Multi-channel carcinogen exposure is superadditive;
   same-channel exposure is additive.}
   Two carcinogens each weakening a \emph{different} coupling channel
   should produce cancer rates exceeding the sum of their individual
   rates. If carcinogen~$X$ degrades $\lambda_i$ and carcinogen~$Y$
   degrades $\lambda_j$ ($i \neq j$), simultaneous exposure
   accelerates the ratchet because each channel loss destabilizes
   the remaining channels (Section~\ref{sec:directionality}): the
   cell loses multiple constraints simultaneously, triggering
   cascading displacement of tissue-referencing propositions in
   $K_0$. Two carcinogens hitting the \emph{same} channel add dose
   to a single degradation process, predicting merely additive risk.
   Partial channel overlap predicts an intermediate
   (sub-multiplicative) interaction.

   Existing epidemiological data confirm this gradient
   (Table~\ref{tab:carcinogen-interactions}).

\begin{table}[ht]
\centering
\small
\caption{Carcinogen interactions stratified by coupling-channel
overlap. RR = relative risk. Interaction type ranges from
super-multiplicative (low channel overlap) through sub-multiplicative
(partial overlap) to additive (same channel). Data from pooled
analyses and meta-analyses as cited.}
\label{tab:carcinogen-interactions}
\begin{tabular}{@{}llllll@{}}
\toprule
Agent 1 & Agent 2 & Cancer & Interaction & Channel overlap \\
\midrule
Smoking & Alcohol & Head/neck &
  Super-mult. &
  Low \\
 & & & RR $\approx 35$ combined & (solvent vs.\ genotoxin) \\
 & & & \cite{Freedman2007} & \\[4pt]
Aflatoxin & Hepatitis B & Liver (HCC) &
  Multiplicative &
  Low \\
 & & & RR $\approx 59$ combined & (DNA damage vs. \\
 & & & vs.\ $3.4 \times 7.3 \approx 25$ & inflammation) \\[4pt]
Smoking & HPV & Cervical &
  Multiplicative &
  Low \\
 & & & RR $\approx 14$--$27$ & (immunosuppression \\
 & & & \cite{Herrero2003} & vs.\ viral oncogene) \\[4pt]
Smoking & Asbestos & Lung &
  Sub-mult. &
  Partial \\
 & & & RR $\approx 28$ combined & (both genotoxic \\
 & & & vs.\ mult.\ $\approx 54$ & + inflammatory) \\
 & & & \cite{Lee2006} & \\[4pt]
Smoking & Radon & Lung &
  Sub-mult. &
  Partial \\
 & & & BEIR VI: less than & (both cause DNA \\
 & & & multiplicative & damage in bronchi) \\[4pt]
PAH mixture & PAH mixture & Various &
  Additive &
  High \\
 & & & EPA sums PEFs & (same CYP1A1 \\
 & & & (dose-addition) & activation pathway) \\
\bottomrule
\end{tabular}
\end{table}

   The gradient is consistent with the framework's prediction:
   agents damaging independent coupling channels
   ($\lambda_i \perp \lambda_j$) produce multiplicative or
   super-multiplicative risk; agents with partial channel overlap
   produce sub-multiplicative risk; agents acting through the same
   channel produce additive risk. This pattern has been observed
   empirically but lacks a unified theoretical explanation.
   The Armitage--Doll multistage model \cite{Armitage1954} captures
   the same structure formally---two agents affecting different
   ``stages'' yield multiplicative risk, while agents affecting the
   same stage yield additive risk---but does not identify what the
   stages \emph{are}. The coupling-channel framework provides the
   identification: stages are coupling channels, and the interaction
   type is determined by whether the agents degrade the same or
   different channels.

   \textbf{The testable refinement:} for any pair of carcinogens
   whose interaction type is currently unknown, the framework
   predicts the interaction from the coupling-channel assignment
   \emph{before} the epidemiological data are collected. This
   converts the prediction from a post-hoc explanation of known
   interactions to a prospective classification tool.

\item \textbf{Sequential monotherapy produces predictable subtype
   transitions.} If a cancer is treated by targeting coupling
   channel~$\lambda_i$, and upon progression is treated by targeting
   channel~$\lambda_j$, the framework predicts a deterministic
   trajectory through molecular subtypes: the tumor severs
   $\lambda_i$ first (producing a phenotype independent of that
   channel), then severs $\lambda_j$ (producing a phenotype
   independent of both). The \emph{order of therapy determines the
   order of channel severance}, and therefore the sequence of
   intermediate subtypes.

   This is testable against existing clinical data. In ER+ breast
   cancer: endocrine therapy first $\to$ ER loss $\to$ if then
   treated with HER2-targeting $\to$ HER2 loss $\to$ triple-negative
   phenotype. The prediction is that this trajectory is
   \emph{reproducible}: the same therapeutic sequence should produce
   the same subtype transitions across patients, because the
   framework dictates which channel severs next. Different
   therapeutic sequences should produce different intermediate
   phenotypes but converge on the same maximally decoupled endpoint.

   The clinical implication is immediate: \textbf{sequential
   monotherapy is a controlled demolition of the escalation chain},
   one channel at a time. Each line of single-agent therapy removes
   one coupling constraint, and the tumor descends one step. This
   provides a structural argument for simultaneous multi-channel
   intervention over sequential single-channel protocols.

\item \textbf{Age-related cancer incidence tracks coupling
   degradation, not mutation accumulation.} Normal aged tissue
   carries remarkably high burdens of cancer-driver mutations without
   developing cancer \cite{Martincorena2015}. Normal esophageal
   epithelium in elderly individuals contains clonal expansions
   carrying \emph{TP53}, \emph{NOTCH1}, and \emph{PIK3CA}
   mutations---textbook oncogenes---covering most of the tissue
   surface, yet cancer does not develop. The framework explains this:
   the mutations are present but the coupling channels are intact.
   The cell carries a \emph{TP53} mutation but still receives
   tissue-level signals, still processes them, still maintains
   tissue-referencing propositions in $K_0$. The mutation weakened
   one constraint but did not sever the escalation chain.

   The framework predicts that cancer incidence across tissues should
   correlate with \textbf{coupling degradation rate}---the rate at
   which tissue-level signaling fidelity declines with age---rather
   than with mutation accumulation rate. This generates a
   distinguishing prediction against SMT: tissues with high mutation
   rates but robust coupling infrastructure (esophagus: high
   turnover, strong architectural signaling) should have
   disproportionately \emph{low} cancer incidence relative to their
   mutational burden, while tissues with lower mutation rates but
   fragile coupling (pancreas: limited coupling redundancy, heavy
   dependence on ductal architecture) should have
   disproportionately \emph{high} incidence. Existing data on
   tissue-specific mutation rates \cite{Martincorena2015} and
   age-specific cancer incidence are sufficient to test this
   prediction.
\end{enumerate}

% ------------------------------------------------------------------
\section{Directionality of Cancer Progression}
\label{sec:directionality}
% ------------------------------------------------------------------

The chain-severance framework makes a strong prediction about the
direction of cancer evolution: \textbf{progression is unidirectional,
from more coupled to less coupled.}

\subsection{The downward ratchet}

A cancer cell operating under partially severed context retains some
coupling channels. But the retention criterion
(\cite{McCarthy2026c}) drives displacement of tissue-referencing
propositions in $K_0$: propositions grounded in coupling channels
that are weakened ($\lambda_j \to 0$) have diminishing marginal
survival value, and are displaced by propositions serving cell-level
optimization. As tissue-referencing propositions leave $K_0$, the
cell's ability to \emph{maintain} its remaining coupling channels
degrades---it is no longer processing the signals that keep those
channels active. Each loss of coupling removes a constraint that
was supporting the next channel. The result is a downward ratchet:
partial severance tends toward complete severance.

This predicts that \textbf{targeted therapy accelerates channel
severance}. Therapy targeting a specific coupling channel
(e.g., endocrine therapy for estrogen-receptor-positive cancer)
creates selection pressure to sever \emph{that specific channel}.
The cancer evolves to eliminate dependence on the targeted channel,
producing a more decoupled, more aggressive phenotype.

Three clinical phenomena confirm this prediction:
\begin{itemize}
\item \textbf{ER+ breast cancer $\to$ triple negative.}
   Endocrine therapy (tamoxifen, aromatase inhibitors) targets the
   $\lambda_3^{\mathrm{endocrine}}$ coupling channel. Under sustained
   endocrine therapy, a subset of ER+ breast cancers lose estrogen
   receptor expression entirely, converting to triple-negative
   phenotype---more aggressive, less differentiated, resistant to the
   original therapy. The cancer severed the channel being targeted.

\item \textbf{Prostate adenocarcinoma $\to$ neuroendocrine small cell.}
   Androgen deprivation therapy targets
   $\lambda_3^{\mathrm{endocrine}}$ (androgen signaling). Under
   sustained deprivation, prostate cancers can undergo lineage
   plasticity, converting to neuroendocrine small cell
   carcinoma---poorly differentiated, maximally aggressive, and
   independent of the hormonal coupling channel that therapy was
   targeting.

\item \textbf{EGFR+ lung adenocarcinoma $\to$ small cell lung cancer.}
   EGFR inhibitors target a growth-factor signaling pathway. Under
   sustained inhibition, a subset of tumors transform to small cell
   lung cancer---a more decoupled phenotype that no longer depends
   on the targeted pathway.
\end{itemize}

In each case, the clinical term is \emph{lineage plasticity} or
\emph{phenotype switching}. The framework provides the structural
explanation: the cancer is not switching randomly. It is descending
the coupling hierarchy, severing whichever channel therapy is
pressuring, and converging on a more completely decoupled phenotype.
The direction is always the same: downward.

\subsection{Why upward progression does not occur}

Can a cell-level cancer spontaneously re-establish tissue-level
coupling and become less aggressive? The framework predicts
\textbf{no}---not without external restoration of the coupling
channel. Re-coupling requires tissue-level signals that the cell
no longer receives (because the channel is severed) and
tissue-referencing propositions that the cell no longer maintains
(because the retention criterion displaced them). The cell cannot
bootstrap its way back up the escalation chain. It would need the
information it has already lost.

This is why differentiation therapy works when it works: it
\emph{externally provides} the tissue-level signals that re-engage
the coupling channel, restoring $\lambda_j$ from outside the cell.
The cell does not re-differentiate spontaneously; it re-differentiates
because the escalation chain is re-established by the therapy.

\subsection{Therapeutic implication: single-channel targeting is
self-defeating}

The directionality result has a direct therapeutic consequence:
\textbf{therapy targeting a single coupling channel accelerates
severance of that channel, producing a more aggressive cancer.}
This is not a side effect but a structural prediction. The cell,
performing descent on $U$ under $C_n^{(0)}$, will optimize away from
dependence on any channel that therapy is pressuring.

Effective therapy must therefore either:
\begin{itemize}
\item Target \emph{multiple} coupling channels simultaneously, so
   that severing one provides no escape (combination immunotherapy +
   differentiation therapy + microenvironment normalization).
\item Target actions that are necessary under $C_n^{(0)}$
   \emph{regardless of which channels remain}, so the cell cannot
   evolve away from them (Section~\ref{sec:therapy}, hybrid strategy).
\item Restore coupling rather than attack the cell, so that the
   selection pressure runs toward re-integration rather than toward
   further decoupling.
\end{itemize}

The clinical trend toward combination therapy, and away from
sequential monotherapy, is consistent with this prediction. The
framework provides the structural reason: sequential monotherapy
is a ratchet that drives the cancer downward through the coupling
hierarchy, one channel at a time.

% ------------------------------------------------------------------
\section{Classification of Therapies by Chain Restoration}
\label{sec:therapy}
% ------------------------------------------------------------------

The framework yields a three-category classification of cancer
therapies based on their structural relationship to the escalation
chain. This classification predicts durability of response.

\begin{enumerate}
\item \textbf{Remove the decoupling cause.} Eliminate the agent that
   is actively degrading coupling channels---before the ratchet
   completes. Works by halting the severance process, allowing
   intact coupling to reassert itself.
\item \textbf{Restore coupling from outside.} Externally re-establish
   the coupling channel that has been severed, returning
   tissue-level actions to the cell's accessible action space.
   Works even after severance because it provides the signals
   and propositions that $K_0$ has lost.
\item \textbf{Impose selection pressure on a decoupled optimizer.}
   Attack or constrain a cell that is already optimizing under
   cell-level $C_n^{(0)}$. The cell's $\gamma > 0$ and
   $\eta_M > 0$ guarantee adaptation; resistance is structurally
   predicted.
\end{enumerate}

\noindent The critical distinction is between \emph{targeting a
coupling channel} (which accelerates severance---the
self-defeating prediction of Section~\ref{sec:directionality})
and \emph{removing whatever is degrading the coupling channel}
(which halts severance). The same signaling pathway can be
involved in both, but the direction of the intervention is
opposite.

\textbf{Molecular subtypes as coupling-channel subtypes.}
Current oncology classifies cancers by molecular markers: ER+,
HER2+, triple-negative breast cancer; luminal vs.\ basal subtypes.
The framework predicts that these molecular subtypes correspond to
\emph{which coupling channel was load-bearing} for that tissue.
ER+ breast cancer is a cancer where
$\lambda_3^{\mathrm{endocrine}}$ was the primary restraint.
Triple-negative breast cancer is a cancer where tissue-architecture
coupling ($\lambda_1$) has also severed. If this mapping holds,
therapy selection should target \emph{restoration of the specific
severed channel}, not merely the expressed receptor. The receptor
is a symptom of which channel severed; the channel is the
therapeutic target.

\subsection{Removing the Decoupling Cause}

These interventions eliminate the agent that is actively degrading
coupling channels. They work when coupling has been weakened but not
yet fully severed---when the cell still retains tissue-referencing
propositions in $K_0$ and the retention-criterion cascade has not
completed.

\textbf{Environmental remediation.}
Removing a subject from a carcinogenic environment---radiation
exposure, asbestos, industrial toxins---halts the ongoing damage to
coupling-channel infrastructure. The coupling channels that have not
yet severed can reassert themselves because the tissue-level signals
are still being sent and the cell still has the propositions needed
to process them. This is why prevention and early removal from
exposure are effective, while removing the original cause after an
established tumor has formed is not: once the ratchet has completed
and tissue-referencing propositions have been displaced from $K_0$,
the cell can no longer re-couple even if the damaging agent is
removed.

\textbf{Pathogen eradication.}
\emph{H.~pylori} eradication reduces gastric cancer incidence
because the bacterium's chronic inflammatory damage is the
decoupling cause for gastric epithelial coupling channels. Hepatitis
B/C treatment reduces hepatocellular carcinoma for the same
structural reason. In each case, removing the pathogen halts the
ongoing degradation of tissue-level signaling.

\textbf{Oncogenic driver inhibition.}
Imatinib in chronic myeloid leukemia (CML) inhibits BCR-ABL, a
constitutively active fusion kinase created by chromosomal
translocation. BCR-ABL is not a coupling channel; it is an aberrant
signal that \emph{overrides} the cell's normal coupling channels,
driving proliferation regardless of tissue-level growth/quiescence
signals. Inhibiting BCR-ABL does not sever a coupling channel---it
removes the signal that was severing them. With the decoupling cause
silenced, the cell's intact coupling channels reassert themselves.

This explains why imatinib produces durable decade-long remissions
rather than the resistance escalation predicted for coupling-channel
targeting. It also explains the nature of imatinib resistance when
it does occur: resistant clones restore the decoupling signal
(T315I mutation, BCR-ABL amplification) rather than descending
further down the coupling hierarchy. The cancer is trying to
re-establish the decoupling cause, not sever additional channels.

\textbf{The timing prediction.} All three examples share a
structural prediction: removing the decoupling cause is effective
\emph{before} the ratchet completes but not after. Asbestos
abatement prevents new mesotheliomas but does not shrink established
tumors. \emph{H.~pylori} eradication prevents gastric cancer but
does not reverse existing malignancy. Imatinib works because
BCR-ABL is the \emph{ongoing} decoupling cause---remove it and the
coupling reasserts. If the CML progressed to blast crisis (fully
decoupled), imatinib response rates drop sharply, consistent with
the ratchet having completed.

\subsection{Chain-Restoring Therapies}

These therapies target the coupling coefficients $\lambda_j$,
externally re-establishing the cell's connection to higher
organizational levels. Unlike removing the decoupling cause, chain
restoration can work even after severance is complete, because it
provides the tissue-level signals from outside the cell.

\textbf{Differentiation therapy (Level $0 \to 1$ restoration).}
All-trans retinoic acid (ATRA) in acute promyelocytic leukemia (APL)
is the paradigm case: it restores the cell's tissue-level identity,
re-engaging the differentiation program that chain severance
disabled. ATRA does not kill the cancer cells; it restores their
coupling to the tissue-level escalation chain, returning tissue-level
actions (differentiation, growth arrest) to the accessible action
space $A_0$. The result is a $>90\%$ cure rate, the highest of any
cancer therapy. On our account, this is not a coincidence but a
\emph{prediction}: therapies that restore the chain should
outperform therapies that do not.

\textbf{Immune checkpoint inhibition (Level $2 \to 3$ restoration).}
PD-1/PD-L1 and CTLA-4 inhibitors restore the organism-level
surveillance coupling $\lambda_3^{\mathrm{immune}}$. Cancer cells
that have downregulated immune recognition markers (severing the
organism-level coupling channel) are re-exposed to immune
surveillance. The therapy does not kill cancer cells directly; it
restores the organizational relationship between the immune system
and the tumor. The durable responses observed in a subset of patients
(complete responses lasting years) are consistent with chain
restoration: once organism-level surveillance is re-established, it
is self-maintaining.

\textbf{Microenvironment normalization (Level $0 \to 1$ and $1 \to 2$
restoration).}
Strategies that restore the tissue microenvironment -- ECM
remodeling, stromal normalization, basement membrane repair --
address the tissue-level coupling infrastructure. Low-dose
anti-angiogenic therapy (vascular normalization rather than vascular
destruction) restores the tissue-organ coupling channel
$\lambda_2^{\mathrm{vascular}}$ by normalizing tumor vasculature,
improving drug delivery and re-establishing metabolic coupling to
the organ level.

\textbf{Gap junction restoration (Level $0 \to 1$ restoration).}
Experimental approaches to restore connexin expression (particularly
Cx43) directly target the highest-bandwidth cell-tissue coupling
channel \cite{Aasen2016}. The framework predicts that restoring GJIC
should suppress malignant behavior even without addressing mutations,
because the coupling restoration returns tissue-level actions to the
accessible action space.

\subsection{Selection Pressure on a Decoupled Optimizer}

These therapies impose selection pressure on a cell that is already
optimizing under cell-level bounded context. The cell's $\gamma > 0$
and $\eta_M > 0$ \cite{McCarthy2026d} guarantee that it will adapt.
Resistance is not a failure of the drug but a structural prediction.

\textbf{Coupling-channel targeting.} Therapies that block a coupling
channel directly---endocrine therapy (ER blockade, androgen
deprivation), single-pathway growth factor inhibitors (EGFR, BRAF)
---fall in this category when the targeted pathway serves as a
tissue-level coupling channel. As Section~\ref{sec:directionality}
predicts, this accelerates severance of that channel, driving
the ratchet downward. The cell, performing descent on $U$, severs
dependence on the targeted channel entirely.

Note the contrast with oncogenic driver inhibition (imatinib, above):
both may target kinase pathways, but the structural direction is
opposite. Driver inhibition removes a \emph{decoupling cause};
coupling-channel blockade removes a \emph{coupling signal}. The
framework predicts different outcomes for structurally different
interventions, even when the molecular target class is similar.

\textbf{Cytoreductive therapies without pathway specificity.}
Chemotherapy and radiation kill or damage cells without targeting
any specific coupling channel. The surviving subpopulation has
undergone selection for resistance within $A_0$---slow-cycling
phenotypes, drug efflux, DNA repair enhancement---none of which
require tissue-level coupling.

\textbf{Surgery.} Removes the decoupled entity entirely. This is
effective for localized disease but does not prevent new severance
events in remaining tissue, and cannot address disseminated disease
(metastasis).

\subsection{The Hybrid Strategy}

The framework predicts that the optimal therapeutic strategy combines
chain restoration with exploitation of structural constraints:
\begin{enumerate}
\item Restore coupling where possible (differentiation therapy,
   immunotherapy, microenvironment normalization).
\item For cells where coupling cannot be restored, target actions
   that are \emph{necessary} under cell-level $C_n^{(0)}$, actions
   the cell cannot optimize away from because they are structurally
   required by the bounded-context constraint. The Warburg effect
   (aerobic glycolysis) is a candidate: if glycolytic metabolism is
   \emph{necessary} under cell-level bounded context (because
   oxidative phosphorylation requires tissue-level oxygen delivery
   infrastructure), then targeting glycolysis imposes a constraint
   the cell cannot evolve around.
\item Remove decoupling causes where identifiable (pathogen
   eradication, environmental remediation, oncogenic driver
   inhibition).
\item Use cytoreductive therapies (chemotherapy, radiation)
   only to reduce tumor burden before chain-restoring
   interventions, not as standalone strategies.
\end{enumerate}

\textbf{Clinical evidence.} The relative success of combination
immunotherapy (chain restoration at level $2 \to 3$) over
single-agent chemotherapy (non-restoring) in melanoma, non-small cell
lung cancer, and renal cell carcinoma is consistent with this
prediction. The curative success of ATRA in APL (chain restoration at
level $0 \to 1$) versus the near-universal resistance to
chemotherapy-only regimens in other acute leukemias provides the
strongest existing evidence.

% ------------------------------------------------------------------
\section{Implications}
\label{sec:implications}
% ------------------------------------------------------------------

\textbf{Cancer is not a disease of cells.} It is a disease of
organizational coupling. The cell itself is functioning correctly;
it is optimizing under its bounded context, exactly as the theorems
predict any entity must. The pathology is in the severed escalation
chain, not in the cell.

\textbf{The ``war on cancer'' metaphor is structurally wrong.} Cancer
cells are not invaders to be destroyed; they are former cooperators
who have lost the organizational infrastructure that made cooperation
possible. Killing them (chemotherapy, radiation) treats the symptom.
Restoring the organizational infrastructure (differentiation therapy,
microenvironment normalization, coupling restoration) treats the
cause. The distinction is between fighting an entity and repairing a
relationship.

\textbf{Prevention over cure.} If cancer is chain severance, then
prevention is coupling maintenance. Factors that maintain tissue
organization (anti-inflammatory interventions, microenvironment
homeostasis, gap junction maintenance) should reduce cancer risk
more effectively than factors that reduce mutational burden,
though the two are related, since mutations are one mechanism of
severance.

\textbf{Peto's paradox.} Large, long-lived organisms do not develop
cancer at rates proportional to their cell count and lifespan
\cite{Peto1975}, a puzzle under the somatic mutation theory,
which predicts that more cell divisions should produce more
mutations and therefore more cancer. The escalation chain framework
offers a structural explanation: larger organisms have more
elaborate tissue organization, stronger intercellular coupling
mechanisms, and more escalation levels. Their
$C_n^{\mathrm{eff}}$ is higher, meaning more severance is required
before cell-level optimization dominates. The additional
organizational infrastructure acts as redundancy in the coupling
channels: multiple mechanisms must fail before the chain severs
completely. Peto's paradox, on this account, is not paradoxical:
more cells but also more coupling yields a roughly constant
severance rate.

\textbf{Benign tumors and non-progression.} Most dysplastic lesions,
adenomas, and benign tumors never become malignant \cite{Nowell1976}.
The framework predicts this: partial decoupling ($0 < \lambda_j < 1$)
does not produce full convergence to the cell-level action space.
A cell with residual tissue coupling retains tissue-referencing
propositions in $K_0$ and does not undergo complete displacement
toward cell-level survival content. The retention criterion
(Paper~3, \cite{McCarthy2026c}) only drives displacement when
tissue-specific propositions have zero marginal value---which
requires $\lambda_j \approx 0$, not merely reduced. Most initiated
cells remain in a partially coupled state indefinitely; their
$C_n^{\mathrm{eff}}$ is reduced but not collapsed to $C_n^{(0)}$.

\textbf{The threshold for progression.} The clinical observation
\cite{Folkman1971} that many cancers persist at low grade for years
before rapid progression is explained by this severance gradient.
Partial decoupling is stable while $\lambda_j$ remains above a
threshold. Progression occurs when $\lambda_j$ crosses that threshold
and cascading decoupling begins: each loss of coupling removes a
constraint that was maintaining the remaining coupling channels.

% ------------------------------------------------------------------
\section{Conclusion}
\label{sec:conclusion}
% ------------------------------------------------------------------

Cancer is not random. Cancer is not ancient. Cancer is what happens
when a cell loses its escalation chain.

This paper has applied the bounded-context framework of Papers~1--4
to cancer biology as a speculative but structurally grounded
proposal. The framework offers what the three major cancer theories
have separately lacked: a formal account of \emph{why} organizational
breakdown would produce the specific, convergent phenotype that it
does. SMT identifies the mechanisms of chain severance (mutations).
The atavistic theory identifies the phenotypic pattern (convergence
on unicellular-like behavior). TOFT identifies the organizational
level (tissue, not cell). Our framework unifies them: mutations
sever the escalation chain (SMT), producing a cell that optimizes
under cell-level bounded context (TOFT), which converges on the same
action space that shaped unicellular life (atavistic pattern),
not by reversion but by structural necessity.

The therapeutic predictions align with where clinical oncology is
already heading: differentiation therapy, immunotherapy, and
microenvironment normalization are gaining ground over purely
cytotoxic approaches. The framework provides a structural reason
for this trend: these therapies restore the escalation chain,
while cytotoxic therapies impose selection pressure on a locally
rational optimizer. The framework generates testable predictions
(Section~\ref{sec:predictions}) that distinguish this account from
existing theories.

The framework's contribution is not empirical but structural: it
offers a target for therapy that is not the cell but the coupling,
not the action but the context that determines which actions are
selected. Whether this structural reasoning maps faithfully onto
cancer biology is an empirical question. The predictions are offered
in that spirit.

\section*{Acknowledgments}

The author used Claude (Anthropic) for drafting assistance, literature review, and \LaTeX{} preparation. All theoretical content, proofs, and arguments are the author's own.

\bibliographystyle{plain}
\begin{thebibliography}{99}

\bibitem{Aasen2016}
T.~Aasen, M.~Mesnil, C.~C. Naus, P.~D. Lampe, and D.~W. Laird.
\newblock Gap junctions and cancer: Communicating for 50 years.
\newblock \emph{Nature Reviews Cancer}, 16(12):775--788, 2016.

\bibitem{Cowan2001}
N.~Cowan.
\newblock The magical number 4 in short-term memory: A reconsideration of
  mental storage capacity.
\newblock \emph{Behavioral and Brain Sciences}, 24(1):87--114, 2001.

\bibitem{Brand1975}
K.~G. Brand, L.~C. Buoen, K.~H. Johnson, and I.~Brand.
\newblock Etiological factors, stages, and the role of the foreign body in
  foreign body tumorigenesis: A review.
\newblock \emph{Cancer Research}, 35(2):279--286, 1975.

\bibitem{Aktipis2015}
C.~A. Aktipis, A.~M. Boddy, G.~Jansen, U.~Hibner, C.~C. Maley, and
  G.~S. Wilkinson.
\newblock Cancer across the tree of life: Cooperation and cheating in
  multicellularity.
\newblock \emph{Philosophical Transactions of the Royal Society B},
  370(1673):20140219, 2015.

\bibitem{Folkman1971}
J.~Folkman, E.~Merler, C.~Abernathy, and G.~Williams.
\newblock Isolation of a tumor factor responsible for angiogenesis.
\newblock \emph{The Journal of Experimental Medicine}, 133(2):275--288, 1971.

\bibitem{Engels2011}
E.~A. Engels, R.~M. Pfeiffer, J.~F. Fraumeni, B.~L. Kasiske,
  A.~K. Israni, J.~J. Snyder, R.~A. Wolfe, N.~P. Goodrich,
  A.~I. Bayakly, G.~L. Clarke, and others.
\newblock Spectrum of cancer risk among {US} solid organ transplant
  recipients.
\newblock \emph{JAMA}, 306(17):1891--1901, 2011.

\bibitem{Davies2011}
P.~C.~W. Davies and C.~H. Lineweaver.
\newblock Cancer tumors as {Metazoa} 1.0: Tapping genes of ancient ancestors.
\newblock \emph{Physical Biology}, 8(1):015001, 2011.

\bibitem{Kawai2008}
T.~Kawai, A.~B. Cosimi, T.~R. Spitzer, N.~Tolkoff-Rubin, M.~Suthanthiran,
  S.~L. Saidman, J.~Shaffer, F.~I. Preffer, R.~Ber, G.~Pascual, and
  D.~H. Sachs.
\newblock {HLA}-mismatched renal transplantation without maintenance
  immunosuppression.
\newblock \emph{New England Journal of Medicine}, 358(4):353--361, 2008.

\bibitem{Hanahan2000}
D.~Hanahan and R.~A. Weinberg.
\newblock The hallmarks of cancer.
\newblock \emph{Cell}, 100(1):57--70, 2000.

\bibitem{Hanahan2011}
D.~Hanahan and R.~A. Weinberg.
\newblock Hallmarks of cancer: The next generation.
\newblock \emph{Cell}, 144(5):646--674, 2011.

\bibitem{Herbst1971}
A.~L. Herbst, H.~Ulfelder, and D.~C. Poskanzer.
\newblock Adenocarcinoma of the vagina: Association of maternal stilbestrol
  therapy with tumor appearance in young women.
\newblock \emph{New England Journal of Medicine}, 284(15):878--881, 1971.

\bibitem{McCarthy2026a}
P.~D. McCarthy.
\newblock Graph structure is necessary for information preservation under
  bounded context.
\newblock Manuscript, 2026.

\bibitem{McCarthy2026b}
P.~D. McCarthy.
\newblock Convergence of control architectures to escalation under bounded context.
\newblock Manuscript, 2026.

\bibitem{McCarthy2026c}
P.~D. McCarthy.
\newblock Knowledge and action are inseparable under survival pressure.
\newblock Manuscript, 2026.

\bibitem{McCarthy2026d}
P.~D. McCarthy.
\newblock Evaluation-driven descent, encoding permanence, and the structural
  invariants of intelligence.
\newblock Manuscript, 2026.

\bibitem{Merlo2006}
L.~M.~F. Merlo, J.~W. Pepper, B.~J. Reid, and C.~C. Maley.
\newblock Cancer as an evolutionary and ecological process.
\newblock \emph{Nature Reviews Cancer}, 6(12):924--935, 2006.

\bibitem{Peto1975}
R.~Peto, F.~J.~C. Roe, P.~N. Lee, L.~Levy, and J.~Clack.
\newblock Cancer and ageing in mice and men.
\newblock \emph{British Journal of Cancer}, 32(4):411--426, 1975.

\bibitem{Nowell1976}
P.~C. Nowell.
\newblock The clonal evolution of tumor cell populations.
\newblock \emph{Science}, 194(4260):23--28, 1976.

\bibitem{Sonnenschein2000}
C.~Sonnenschein and A.~M. Soto.
\newblock Somatic mutation theory of carcinogenesis: Why it should be dropped
  and replaced.
\newblock \emph{Molecular Carcinogenesis}, 29(4):205--211, 2000.

\bibitem{Trigos2017}
A.~S. Trigos, R.~B. Pearson, A.~T. Papenfuss, and D.~L. Goode.
\newblock Altered interactions between unicellular and multicellular genes
  drive hallmarks of transformation in a diverse range of solid tumors.
\newblock \emph{Proceedings of the National Academy of Sciences},
  114(24):6406--6411, 2017.

\bibitem{Vogelstein2013}
B.~Vogelstein, N.~Papadopoulos, V.~E. Velculescu, S.~Zhou, L.~A. Diaz, and
  K.~W. Kinzler.
\newblock Cancer genome landscapes.
\newblock \emph{Science}, 339(6127):1546--1558, 2013.

\bibitem{Huang2009}
S.~Huang, I.~Ernberg, and S.~Kauffman.
\newblock Cancer attractors: a systems view of tumors from a gene network
  dynamics and developmental perspective.
\newblock \emph{Seminars in Cell \& Developmental Biology}, 20(7):869--876,
  2009.

\bibitem{Massague2016}
J.~Massagu{\'e} and A.~C. Obenauf.
\newblock Metastatic colonization by circulating tumour cells.
\newblock \emph{Nature}, 529(7586):298--306, 2016.

\bibitem{Warburg1956}
O.~Warburg.
\newblock On the origin of cancer cells.
\newblock \emph{Science}, 123(3191):309--314, 1956.

\bibitem{Zhao2017}
C.-M. Zhao, Y.~Hayakawa, Y.~Kodama, S.~Muthupalani, C.~B. Westphalen,
  G.~T. Andersen, A.~Flatberg, H.~Johannessen, R.~A. Friedman,
  B.~W. Renz, and others.
\newblock Denervation suppresses gastric tumorigenesis.
\newblock \emph{Science Translational Medicine}, 6(250):250ra115, 2014.

\bibitem{Armitage1954}
P.~Armitage and R.~Doll.
\newblock The age distribution of cancer and a multi-stage theory of
  carcinogenesis.
\newblock \emph{British Journal of Cancer}, 8(1):1--12, 1954.

\bibitem{Freedman2007}
N.~D. Freedman, C.~C. Abnet, T.~M. Leitzmann, T.~Mouw, A.~F. Subar,
  A.~R. Hollenbeck, and A.~Schatzkin.
\newblock A prospective study of tobacco, alcohol, and the risk of
  esophageal and gastric cancer subtypes.
\newblock \emph{American Journal of Epidemiology}, 165(12):1424--1433, 2007.

\bibitem{Herrero2003}
R.~Herrero, X.~Castellsagu{\'e}, M.~Pawlita, J.~Lissowska,
  F.~Kee, P.~Balaram, T.~Rajkumar, H.~Sridhar, B.~Rose,
  J.~Pintos, and others.
\newblock Human papillomavirus and oral cancer: the International
  Agency for Research on Cancer multicenter study.
\newblock \emph{Journal of the National Cancer Institute},
  95(23):1772--1783, 2003.

\bibitem{Lee2006}
P.~N. Lee.
\newblock Relation between exposure to asbestos and smoking jointly
  and the risk of lung cancer.
\newblock \emph{Occupational and Environmental Medicine},
  58(3):145--153, 2001.

\bibitem{Martincorena2015}
I.~Martincorena, A.~Roshan, M.~Gerstung, P.~Ellis, P.~Van~Loo,
  S.~McLaren, D.~C. Wedge, A.~Fullam, L.~B. Alexandrov, J.~M. Tubio,
  and others.
\newblock High burden and pervasive positive selection of somatic
  mutations in normal human skin.
\newblock \emph{Science}, 348(6237):880--886, 2015.

\end{thebibliography}

\end{document}
